\documentclass[conference]{IEEEtran}
\usepackage{cite}
\usepackage{amsmath,amssymb,amsfonts}
\usepackage{array}
\usepackage{booktabs}
\usepackage{xcolor}
\usepackage{url}
\newcommand{\pec}{\textsc{PEC}}
\newcommand{\qsl}{\textsc{QSL}}
\newcommand{\qec}{\textsc{QEC}}

\begin{document}

\title{Is There an Upper Limit to Computer Processing Speed?\\
An Evidence-Bounded Research Proposal for Physical Compute Envelopes}

\author{\IEEEauthorblockN{Anonymous Research Proposal}
\IEEEauthorblockA{\textit{Survey to Idea to ExperimentDesign to Author Workflow}\\
IEEE Conference Format Research Plan}}

\maketitle

\begin{abstract}
There are physical limits to computer processing speed, but no single
technology-independent number answers the question for every computer and
every workload. A processor clock, transistor count, physical gate rate,
sampling rate, and application time-to-solution have different denominators.
They cannot be substituted for each other without an explicit computation
boundary. This paper proposes the Physical Compute-Envelope Contract, a
research protocol for evaluating a useful-throughput claim under a declared
workload, correctness target, reliability target, architecture, energy,
thermal environment, spatial extent, communication topology, memory
assumptions, and input-output boundary. The proposal combines quantum speed
limits, finite signal propagation, thermodynamic erasure cost, heat removal,
information storage, and fault-tolerant logical overhead as separate
constraint cards. It rejects two common overclaims: that the end of Moore
scaling supplies a universal physical ceiling, and that quantum computation
removes physical speed limits. Quantum processors can change algorithmic
complexity for specified tasks, but state preparation, control, communication,
measurement, error correction, and verification remain part of useful
end-to-end computation. The resulting design is a falsifiable and
architecture-neutral way to decide whether a proposed speed statement is
supported within scope, incomplete, overgeneralized, or inconclusive. This is
a design-only study: it reports no executed hardware experiment, simulation,
benchmark, or observed performance result.
\end{abstract}

\begin{IEEEkeywords}
physical limits of computation, quantum speed limit, Moore's Law, Landauer
principle, fault-tolerant quantum computing, computer architecture,
throughput, reproducibility
\end{IEEEkeywords}

\section{Introduction}

The statement that computers will eventually become as fast as physics allows
is compelling but underspecified. What counts as "fast"? A contemporary
processor may advertise a clock frequency, an accelerator may advertise
operations per second, a high-performance system may advertise floating-point
throughput, and a quantum processor may report physical-gate performance,
sampling throughput, or a task-specific computational advantage. None of
these numbers independently specifies the useful work delivered to an end
user. The computation could differ in problem size, output representation,
accuracy, probability of success, initialization, readout, energy boundary,
or communication path.

The historical motivation is clear. Moore's influential observation concerned
the number of components that could be integrated economically over time
\cite{moore}. Earlier idealized scaling also tied smaller field-effect devices
to favorable density, speed, and power trends \cite{dennard}. Those
relationships enabled extraordinary engineering progress, but they are not
laws that a physical computer must obey. Leakage, variability, interconnects,
power delivery, cooling, packaging, manufacturing cost, and architecture
change the realized trajectory. The reported slowing of Moore-style scaling is
therefore a reason to ask a physics question, not an answer to it
\cite{waldrop}.

Quantum information processing sharpens rather than removes the question.
Quantum speed limits constrain how quickly specified quantum states can
evolve under stated physical resources. Quantum algorithms can outperform
classical alternatives for particular problem structures. Yet neither result
creates a general-purpose, unbounded operations-per-second machine. A useful
computation remains embodied in a device with finite energy, finite size,
finite signal velocity, control hardware, thermal interfaces, memory,
measurement, and a required probability of correct output.

This report simulates the repository workflow
\begin{equation}
\text{Survey} \longrightarrow \text{Idea} \longrightarrow
\text{ExperimentDesign} \longrightarrow \text{Author}.
\label{eq:workflow}
\end{equation}
It gives a direct answer: \emph{there are upper limits, but they are
conditional physical compute envelopes rather than one universal speed
constant}. The proposed contribution is not a new numerical limit. It is a
testable method for stating which limit is active for one clearly defined
claim.

\section{Survey: Why a Single Speed Number Fails}

\subsection{A metric has a numerator, denominator, and scope}

Table~\ref{tab:metrics} separates quantities often called computer speed.
They should be compared only after their numerator, time boundary, and
acceptance rule have been aligned. A clock is a periodic control signal; it is
not an accepted answer rate. A physical quantum gate rate is not a logical
algorithm rate. A sampling rate is not necessarily a rate for a practically
useful task. Conversely, a slow physical operation may contribute to a faster
algorithm if it changes computational complexity for the named workload.

\begin{table}[t]
\caption{Common Speed Metrics and Their Missing Context}
\label{tab:metrics}
\centering
\footnotesize
\begin{tabular}{p{0.20\columnwidth}p{0.29\columnwidth}p{0.29\columnwidth}}
\toprule
Metric & What it records & What it does not establish\\
\midrule
Clock frequency & Local timing cadence & Useful work, memory behavior, energy, or I/O cost\\
Transistor count & Integration density & Architecture, yield, heat, or algorithmic work\\
Physical gate rate & Device-level operation time & Logical reliability and full algorithm completion\\
FLOPS or TOPS & Chosen arithmetic operation rate & Accuracy, data movement, and workload relevance\\
Sampling rate & Samples from a named distribution & General-purpose advantage or application utility\\
Time-to-solution & End result for one workload & Transfer to a different workload or condition\\
\bottomrule
\end{tabular}
\end{table}

The scope problem can be made explicit. Let \(W\) denote a workload including
its inputs and output representation, \(\epsilon\) a correctness or
distributional-error target, and \(\rho\) a reliability target. A claim about
useful throughput is meaningful only after it states an end-to-end time
boundary. The direct dependent quantity in this proposal is

\begin{equation}
R_{\mathrm{accepted}} =
\frac{N_{\mathrm{accepted}}}{T_{\mathrm{wall}}},
\label{eq:accepted-rate}
\end{equation}

where \(N_{\mathrm{accepted}}\) is the number of outputs meeting the declared
\(\epsilon\) and \(\rho\) criteria and \(T_{\mathrm{wall}}\) includes or
explicitly excludes preparation, communication, readout, verification, and
classical postprocessing. Equation~\eqref{eq:accepted-rate} is a definition,
not a claim of measured performance.

\subsection{Moore scaling is an engineering trajectory, not a terminal theorem}

Moore's original article was a projection about component density and
cost-effective integration \cite{moore}. Dennard and colleagues analyzed a
device-scaling regime in which electric fields and device dimensions were
coordinated \cite{dennard}. The importance of these sources is conceptual:
they show why density, power, and speed were historically coupled. Their
importance is not that either predicts an unavoidable absolute speed for
future computers.

At small feature sizes, multiple bottlenecks decouple. A switching device may
be fast while a wire, cache miss, package, analog interface, or thermal
constraint determines wall-clock completion. Three-dimensional integration,
special-purpose acceleration, better algorithms, reversible logic, optical
interconnects, and heterogeneous systems can change the active bottleneck.
The appropriate inference from an atomic-scale feature constraint is that
specific integration routes have limits; it is not that every information
processing architecture has reached one scalar ceiling.

\subsection{Physical bounds constrain different mechanisms}

Quantum speed limits offer a strong example of a real physical constraint that
is often overextended. For an ideal closed quantum system evolving to an
orthogonal state, the Mandelstam-Tamm and Margolus-Levitin forms can be
written as

\begin{equation}
\tau_{\mathrm{QSL}} \geq
\max\left\{\frac{\pi\hbar}{2\Delta E},
\frac{\pi\hbar}{2(E-E_0)}\right\},
\label{eq:qsl}
\end{equation}

where \(\Delta E\) is the energy standard deviation, \(E-E_0\) is mean energy
above the ground-state reference, and \(\hbar\) is reduced Planck's constant.
The review by Deffner and Campbell organizes these bounds as limits on
quantum evolution and control \cite{deffner}; the original
Margolus-Levitin result provides a related dynamical speed statement
\cite{margolus}. Lloyd's physical-computation analysis provides a wider
resource-bound perspective without making energy the sole whole-system
denominator \cite{lloyd}. Equation~\eqref{eq:qsl} says neither that every quantum gate
saturates the bound nor that the reciprocal of this time is a complete
computer rate.

Finite transmission produces a different lower bound. If the critical
information path has length \(L\) and signal velocity
\(v_{\mathrm{sig}}\), then

\begin{equation}
\tau_{\mathrm{signal}} \geq \frac{L}{v_{\mathrm{sig}}},
\qquad v_{\mathrm{sig}}\leq c,
\label{eq:signal}
\end{equation}

with \(c\) the speed of light in vacuum. The relevant velocity can be lower
in a medium, and a real path also contains serialization, switching, and
controller delay. This term applies to a named topology; it is not a
replacement for a device-transition model.

Landauer's principle addresses logically irreversible erasure. At temperature
\(T\), erasing \(n_{\mathrm{erase}}\) independent bits has an ideal lower heat
cost

\begin{equation}
Q_{\mathrm{erase}} \geq n_{\mathrm{erase}} k_{\mathrm{B}}T\ln 2,
\label{eq:landauer}
\end{equation}

where \(k_{\mathrm{B}}\) is Boltzmann's constant \cite{landauer}. The
assumptions matter: \eqref{eq:landauer} is not a charge for every logical
operation and not a measured package power. It is nevertheless indispensable
when a claim relies on repeated reset, irreversible control, measurement, or
data disposal. Finite-time operations and heat extraction can add substantial
cost above this ideal floor.

\subsection{Quantum advantage requires task semantics}

Quantum information can deliver algorithmic advantages, but the predicate is
not "uses qubits." Preskill describes the central role of noise and circuit
depth for near-term quantum systems \cite{preskill}, while DiVincenzo
emphasizes the physical requirements for initialization, control,
measurement, and communication \cite{divincenzo}. Random-sampling studies
clarify that a computational-advantage claim depends on a task distribution,
classical comparator, verification method, and complexity assumptions
\cite{hangleiter}.

Photonic boson-sampling demonstrations are scientifically important
task-specific evidence. They should not be reported as a universal rate at
which arbitrary programs become faster. A particular classical-runtime
estimate can change with improved algorithms and hardware; the sampling task
may have no direct application utility; and a universal, fault-tolerant
machine has additional preparation, error-handling, and readout requirements.
The survey therefore retains sampling as an architecture and workload card,
not as proof that an upper speed limit vanished.

\section{Idea: The Physical Compute-Envelope Contract}

\subsection{From a broad question to a falsifiable object}

The Idea Agent selected the \pec{} after rejecting a clock-only comparison,
an energy-only ceiling, and a quantum escape narrative. The research object is
one tuple

\begin{equation}
\mathcal{C}=\langle W,\epsilon,\rho,A,E,T,L,\mathcal{N},
\mathcal{H},\mathcal{M},\mathcal{I}\rangle,
\label{eq:contract}
\end{equation}

where \(A\) identifies hardware and algorithmic architecture;
\(\mathcal{N}\) identifies communication topology; \(\mathcal{H}\) identifies
the cooling and heat-extraction boundary; \(\mathcal{M}\) identifies memory
and information-density assumptions; and \(\mathcal{I}\) identifies
input-output, observation, and verification boundaries. The other symbols are
defined in Section II. A missing field is not a nuisance parameter; it is an
unresolved part of the speed claim.

The PEC central hypothesis is that a layered contract distinguishes supported
from unsupported useful-throughput statements more reliably than one component
metric. It can be falsified. The framework fails if its complete card set
cannot discriminate a well-supported claim from an unsupported claim, if a
mandatory card adds no decision value across the declared workload family, or
if independently verified logical throughput violates the stated envelope
without a changed contract or accounting correction.

\begin{figure*}[t]
\centering
\fbox{\begin{minipage}{0.92\textwidth}
\centering
\vspace{4pt}
\textbf{Claim Contract}\\[2pt]
\(\downarrow\)\\[-2pt]
\textbf{Transition Card} \quad
\(\textbf{Signal and I/O Card}\) \quad
\textbf{Thermal and Memory Card} \quad
\textbf{Logical Reliability Card}\\[2pt]
\(\downarrow\)\\[-2pt]
\textbf{PEC Envelope: applicability, exclusions, active constraints}\\[2pt]
\(\downarrow\)\\[-2pt]
\textbf{Conditional Decision: in-scope support, hold, redesign, or inconclusive}
\vspace{4pt}
\end{minipage}}
\caption{PEC claim-to-decision architecture. The delivered editable Draw.io
source expresses the same relationship. The diagram is a design artifact, not
an executed benchmark.}
\label{fig:pec-flow}
\end{figure*}

\subsection{Compositional envelope rather than a false universal theorem}

For a given contract, PEC defines the following reporting envelope:

\begin{align}
R_{\mathrm{accepted}} \leq
\min\{R_{\mathrm{transition}},R_{\mathrm{signal}},
R_{\mathrm{thermal}}, \nonumber\\
R_{\mathrm{memory}},R_{\mathrm{reliable}},R_{\mathrm{I/O}}\}.
\label{eq:pec-envelope}
\end{align}

Equation~\eqref{eq:pec-envelope} is deliberately operational. It does not
assert that the six right-hand terms are independent, tight, or exhaustive
physical theorems. Instead, it prevents the more common error of reporting a
favorable numerator while silently omitting a limiting system stage. A card
may be marked inapplicable with a documented rationale; it may not be
silently assumed favorable.

Table~\ref{tab:cards} specifies the minimum information in each card. For
example, a reported quantum physical-gate time belongs in the transition card.
It cannot establish logical useful throughput until the reliability card
states encoding, error model, syndrome cycles, decoder, retries, and
accepted-output criterion. A high-speed chip-to-chip link belongs in the
signal card. It cannot erase I/O cost if data preparation or output
verification crosses a slower boundary.

\begin{table*}[t]
\caption{PEC Cards Required Before a Useful-Throughput Claim}
\label{tab:cards}
\centering
\footnotesize
\begin{tabular}{p{0.19\textwidth}p{0.36\textwidth}p{0.35\textwidth}}
\toprule
Card & Minimum evidence record & Claim invalidated when absent\\
\midrule
Workload semantics & Input distribution, output form, accuracy metric, reliability threshold, wall-clock start and stop. & Different computations are treated as the same unit of work.\\
Transition & Switching or state-evolution model, control assumptions, and QSL applicability rationale. & A device-local time is promoted to a system rate.\\
Signal and I/O & Critical paths, topology, media, controller paths, initialization, readout, and verification inclusion. & Data movement or observation is silently excluded.\\
Thermal and memory & Irreversible operations, temperature, heat-removal limit, density assumptions, and finite-time losses. & Ideal erasure cost is reported as realized dissipation.\\
Logical reliability & Physical error model, code, space-time overhead, decoder, postselection or retry rule, and output acceptance probability. & Physical operations are counted as logical answers.\\
Transfer and comparator & Classical comparator, resource normalization, architecture version, and conditions of transfer. & A task-specific result is framed as universal speed.\\
\bottomrule
\end{tabular}
\end{table*}

\section{ExperimentDesign: A Design-Only Evaluation Protocol}

\subsection{Research unit and variable map}

The ExperimentDesign Agent routes PEC to a computational-digital and
mathematics-theory template. The unit of analysis is

\begin{equation}
U=\langle W,\epsilon,\rho,A,
\tau_{\mathrm{transition}},\tau_{\mathrm{signal}},
\mathcal{H},\mathcal{M},\Omega_{\mathrm{logical}},\mathcal{I}\rangle,
\label{eq:unit}
\end{equation}

where \(\Omega_{\mathrm{logical}}\) is a physical-to-logical space-time and
control-overhead description. The independent variables are architecture class
and workload class. The dependent variable is
\(R_{\mathrm{accepted}}\) from \eqref{eq:accepted-rate}. Controlled
conditions include the correctness definition, reliability target, output
representation, and time boundary.

The protocol cannot report a scalar readiness score. A scalar would hide
whether an attractive transition time was paid for by an implausible cooling
assumption, an omitted decoder, a changed output distribution, or an
unreported classical preprocessing step. Instead it provides traceable cards
and conditional decisions.

\subsection{Pre-registered reasoning obligations}

First, every transition claim must say whether \eqref{eq:qsl} applies. A QSL
may be inapplicable to an open, driven, nonorthogonal, or insufficiently
specified process. Inapplicability is a valid result of the card, but
substituting a reciprocal gate time for the QSL is not.

Second, the transport card must list the longest critical information path and
the relevant medium. Equation~\eqref{eq:signal} is only a propagation floor;
queueing, routing, synchronization, and decoder feedback can be additional
terms. A small local quantum processor can therefore be control-limited,
while a large conventional processor can be memory- or network-limited.

Third, the thermal card separates a fundamental lower bound from a realized
engineering budget. Equation~\eqref{eq:landauer} applies to declared
irreversible erasures. A future study would document which resets,
measurements, and control records count, then compare their dissipation and
heat removal with the package boundary. Reversible logic is a crucial
counterexample: it can change the erasure accounting, but it does not make
signal latency, finite device volume, or error handling disappear.

Fourth, the reliability card converts physical process counts into accepted
logical work. A useful reporting relation is

\begin{equation}
R_{\mathrm{logical}} =
\frac{R_{\mathrm{physical}}Y_{\mathrm{accept}}}
{\Omega_{\mathrm{space}}\Omega_{\mathrm{time}}\Omega_{\mathrm{control}}},
\label{eq:logical-rate}
\end{equation}

where \(Y_{\mathrm{accept}}\) is the probability that a result passes the
declared acceptance rule, and the three \(\Omega\) terms are documented
resource multipliers. Equation~\eqref{eq:logical-rate} is a bookkeeping
model, not a claim that all architectures factor exactly this way. It exposes
why raw physical-gate speed is insufficient for a reliable algorithm.

\subsection{Future evidence and conditional branches}

No experiment, simulation, device operation, or cloud job is performed in
this report. Table~\ref{tab:future} states what a future research program
could test after qualified human review. It also prevents each evidence source
from carrying a claim it cannot establish.

\begin{table}[t]
\caption{Future Evidence Modalities and Their Boundaries}
\label{tab:future}
\centering
\footnotesize
\begin{tabular}{p{0.27\columnwidth}p{0.25\columnwidth}p{0.26\columnwidth}}
\toprule
Evidence modality & Can test & Cannot establish alone\\
\midrule
Source-reviewed device record & Populated transition, topology, and environment fields & End-to-end accepted throughput\\
Fixed-workload benchmark & Scoped correctness and time boundary & Universal physical ceiling\\
QEC resource analysis & Logical overhead under one error model & Performance under another hardware model\\
Thermal/package assessment & Compatibility with a named heat boundary & Minimum cost of all computation\\
Independent replication & Transfer under the same contract & Transfer to an unlisted workload or scale\\
\bottomrule
\end{tabular}
\end{table}

The protocol yields six statuses. CONTRACT EVALUABLE means that all cards and
exclusions are explicit, not that a speed claim is true. BOUND SUPPORTED IN
SCOPE means that future evidence supports the named envelope only under the
declared contract. RELIABILITY ACCOUNT INCOMPLETE holds a claim when physical
operations have not been translated into accepted logical outputs.
THERMAL OR I/O BOUNDARY MISSING requires a redesign when cooling, preparation,
readout, or verification has been omitted. QUANTUM ADVANTAGE OVERGENERALIZED
requires the task and comparator to be separated from a general processor
claim. INCONCLUSIVE is an allowed outcome when evidence cannot discriminate
the envelope.

\section{What Quantum Computing Does and Does Not Change}

Quantum computation matters because it can alter the relationship between
workload size and required computational resources. It does not matter because
it removes the requirement for a resource boundary. The DiVincenzo criteria
remain a useful implementation checklist: scalable qubits, initialization,
long coherence relative to operations, a universal gate set, measurement, and
communications are physical requirements, not abstract algorithm steps
\cite{divincenzo}.

Error correction makes the distinction sharper. Surface-code research
describes a practical route to large-scale fault tolerance while emphasizing
the associated architecture \cite{fowler}. More recent cross-validated work
reports low-overhead fault-tolerant memory and alternative code families
\cite{bravyi,cohen}. These sources justify optimism about reducing overhead;
they do not license the omission of overhead. The relevant performance object
is a logical computation with a target failure probability, not a raw
physical-qubit cycle.

Random sampling demonstrates a related lesson. The photonic experiment by
Zhong and colleagues is an important milestone in quantum computational
advantage for a defined sampling task \cite{zhong}. The subsequent review
literature explains why the task, verification procedure, and classical
simulation assumptions are part of the meaning of such a claim
\cite{hangleiter}. PEC records those conditions. It does not discount the
demonstration; it blocks an invalid transfer from a specialized sampling
experiment to "quantum computers are unconstrained by processing speed."

\section{Expected Outcome Branches and Scientific Payoff}

The anticipated contribution is a decision discipline, not a forecast.
Table~\ref{tab:outcomes} pre-registers allowable interpretations. It makes
clear that an evidence-complete card set can authorize specialist evaluation
without declaring deployment readiness, and that a favorable bounded result
does not predict a universal successor technology.

\begin{table*}[t]
\caption{Pre-Registered Conditional Interpretations}
\label{tab:outcomes}
\centering
\footnotesize
\begin{tabular}{p{0.25\textwidth}p{0.36\textwidth}p{0.29\textwidth}}
\toprule
Status & Authorized future interpretation & Required action\\
\midrule
CONTRACT EVALUABLE & The claim is sufficiently specified for a physical and architecture review. & Do not call it a measured speed or deployment approval.\\
BOUND SUPPORTED IN SCOPE & A named active constraint is consistent with future evidence for the named workload and architecture. & Preserve all exclusions and compare only compatible contracts.\\
RELIABILITY ACCOUNT INCOMPLETE & Raw operations cannot yet be interpreted as accepted logical work. & Add an error model, code, overhead, and acceptance criterion.\\
THERMAL OR I/O BOUNDARY MISSING & The system boundary omits a possible limiting process. & Add cooling, data movement, preparation, readout, or verification evidence.\\
QUANTUM ADVANTAGE OVERGENERALIZED & A specialized computational demonstration is being transferred beyond its task. & State the task, comparator, verification method, and missing utility bridge.\\
INCONCLUSIVE & The evidence does not select an active envelope term. & Report no direction for a numerical upper limit.\\
\bottomrule
\end{tabular}
\end{table*}

The scientific payoff is comparability. A conventional accelerator and a
quantum processor need not have identical mechanisms to be compared fairly.
They need compatible workload contracts. If their answers have different
accuracy, success probability, I/O inclusion, energy boundary, or latency
boundary, the correct result is not a leaderboard. It is a statement that
the comparison is underdetermined.

PEC also creates useful negative knowledge. A claim may be technically
impressive but nontransferable: a new gate may reduce
\(\tau_{\mathrm{transition}}\) without changing
\(\tau_{\mathrm{signal}}\); a better code may reduce
\(\Omega_{\mathrm{space}}\) while increasing controller complexity; a
sampling benchmark may establish a nonclassical distribution without
establishing application time-to-solution. These are not failures. They
identify the specific next research dependency.

\section{Risks, Limitations, and Human Review}

The proposal has deliberate limits. First, a minimum-of-terms envelope can
miss couplings between heat, error rate, control bandwidth, and architecture.
The purpose of \eqref{eq:pec-envelope} is disclosure, not a proof that every
real system is separable. A future implementation must model couplings or
flag them as unknown.

Second, physical bounds depend on assumptions. The energy reference in a
Margolus-Levitin argument, the state distance in a QSL, the temperature and
logical irreversibility in a Landauer calculation, and the signal path in a
latency analysis all require careful source review. A numerical value should
never be inserted merely because a formula exists.

Third, algorithmic complexity and classical simulation methods evolve.
Consequently, a claim of quantum advantage needs a stated classical comparator
and date-specific method, not only a historical runtime estimate. A different
classical algorithm can change the benchmark meaning without changing the
quantum device.

Finally, qualified device physicists, computer architects,
quantum-information researchers, thermal engineers, and reproducibility
reviewers must assess any real parameterization. They must verify source
applicability, units, device version, error model, workload specification,
software stack, measurement inclusion, and independent replication protocol.
This report is not a safety certification, a hardware control procedure, or
an empirical validation.

\section{Research Priorities Enabled by PEC}

\subsection{A three-family validation suite}

The immediate research priority is not to search for a grand universal
constant. It is to make upper-limit arguments comparable across three
deliberately different workload families. The first family should be a
conventional memory-intensive workload, for which data movement, cache
hierarchy, interconnect, packaging, and heat extraction may dominate a local
arithmetic rate. The second should be a highly structured numerical workload,
for which a specialized accelerator may improve useful throughput even if its
clock does not. The third should be a quantum workload with a formally stated
input distribution, output verification criterion, and classical comparator.
This suite prevents the protocol from being tuned to one hardware culture.

For each family, the baseline and candidate must share a stated output
semantic. A conventional method returning an approximate scalar estimate
cannot be compared directly with a quantum sampler returning a distribution
unless both are mapped to an agreed accuracy and utility criterion. A future
study must therefore publish the workload generator, instance class, error
metric, reliability threshold, and end-to-end stop condition before comparing
rates. PEC treats this publication requirement as a scientific control, not
as administrative documentation.

\subsection{Separate fundamental floors from engineering ceilings}

The proposed program also separates three kinds of statement that are often
mixed together. A fundamental floor is an ideal constraint derived under
explicit physical assumptions, such as \eqref{eq:qsl} or
\eqref{eq:landauer}. An engineering ceiling is a system-specific limit set by
a package, controller, interconnect, fabrication process, or cooling system.
A benchmark outcome is a performance observation on one workload and one
software stack. These layers can inform each other, but none should be
reported as the other.

This distinction gives direct hypotheses to future experiments. If a device
improvement lowers an observed time-to-solution, the resource cards can ask
whether the improvement came from a faster transition, a shorter data path, a
better compiler schedule, lower error-correction overhead, changed
postselection, relaxed accuracy, or a changed time boundary. If the answer is
not known, the result remains valuable but cannot be used to update a
fundamental upper-limit claim. Conversely, if an engineering ceiling becomes
active far above an ideal floor, research may target the named engineering
constraint rather than portraying the gap as a failure of physics.

\subsection{A claim-reporting contract for reproducible comparisons}

Every future PEC report should contain a compact machine-readable contract in
addition to prose. The required fields are the workload identifier, output
type, input distribution, accuracy target, reliability target, processor and
compiler version, physical topology, timing boundary, energy and thermal
boundary, memory accounting, communication paths, logical-overhead model,
classical comparator, and uncertainty interval. Each field should record one
of four states: supported by a primary source, independently verified in the
future study, inapplicable with rationale, or unknown.

The contract is intentionally strict about unknowns. A missing cooling
specification does not become an implicit infinite heat sink. An omitted
decoder does not become zero latency. A missing classical preprocessing step
does not become free. This rule is particularly important for cross-platform
claims, because the largest apparent speedup can result from moving a cost
outside the denominator rather than from changing the computation itself.

\subsection{Predictions that can be wrong}

PEC produces several falsifiable predictions about research reporting. First,
after I/O, reliability, and thermal cards are added, some raw device-rate
rankings will change for useful end-to-end workloads. Second, the active
constraint will often change as an architecture scales: a local transition
term may dominate a small prototype while wiring, synchronization, error
correction, or cooling dominates a larger system. Third, claims with the
largest transfer distance from a specialized benchmark to a general-purpose
conclusion will show the greatest number of unknown cards.

These predictions can be tested by independent analysts using public
specifications, controlled workloads, and documented assumptions. A negative
result would be informative. If PEC cards never change an interpretation or
if their decision statuses disagree systematically with independently
validated comparisons, the framework should be simplified or rejected. The
proposal therefore makes its own methodology subject to the same evidence
discipline it asks of computing-speed claims.

\section{Conclusion}

Computer processing speed has physical upper limits, but the correct
scientific answer is not one scalar measured in hertz, FLOPS, gates per
second, or transistor count. A valid limit is conditional on what the
computer is asked to compute, how correctness and reliability are defined,
where wall-clock time begins and ends, and which physical resources are
allowed. Moore-style scaling remains historically important but is not a law
of physics. Quantum computing can alter algorithmic resource scaling for
specified tasks but does not bypass state-evolution, signal, thermodynamic,
thermal, information, reliability, control, and I/O constraints.

The Physical Compute-Envelope Contract converts that conclusion into a
falsifiable research program. It requires a workload contract and six linked
constraint cards before a useful-throughput claim can be evaluated. It thereby
replaces two misleading narratives - a universal terminal silicon speed and
an unconstrained quantum future - with a more direct research question:
which constraint controls a clearly stated computation, and what evidence
would show that the constraint has changed? This is a stronger basis for
comparing architectures and for designing the experiments that must follow.

\begin{thebibliography}{00}
\bibitem{moore}
G. E. Moore, "Cramming more components onto integrated circuits,"
\emph{Electronics}, vol. 38, no. 8, pp. 114--117, 1965.

\bibitem{dennard}
R. H. Dennard, F. H. Gaensslen, H.-N. Yu, V. L. Rideout, E. Bassous, and
A. R. LeBlanc, "Design of ion-implanted MOSFET's with very small physical
dimensions," \emph{IEEE J. Solid-State Circuits}, vol. 9, no. 5,
pp. 256--268, 1974, doi: 10.1109/JSSC.1974.1050511.

\bibitem{waldrop}
M. M. Waldrop, "The chips are down for Moore's law," \emph{Nature},
vol. 530, pp. 144--147, 2016, doi: 10.1038/530144a.

\bibitem{margolus}
N. Margolus and L. B. Levitin, "The maximum speed of dynamical evolution,"
\emph{Physica D}, vol. 120, no. 1--2, pp. 188--195, 1998,
doi: 10.1016/S0167-2789(98)00054-2.

\bibitem{deffner}
S. Deffner and S. Campbell, "Quantum speed limits: from Heisenberg's
uncertainty principle to optimal quantum control," \emph{J. Phys. A:
Math. Theor.}, vol. 50, no. 45, Art. no. 453001, 2017,
doi: 10.1088/1751-8121/aa86c6.

\bibitem{landauer}
R. Landauer, "Irreversibility and heat generation in the computing process,"
\emph{IBM J. Res. Dev.}, vol. 5, no. 3, pp. 183--191, 1961,
doi: 10.1147/rd.53.0183.

\bibitem{lloyd}
S. Lloyd, "Ultimate physical limits to computation," \emph{Nature},
vol. 406, pp. 1047--1054, 2000, doi: 10.1038/35023282.

\bibitem{divincenzo}
D. P. DiVincenzo, "The physical implementation of quantum computation,"
\emph{Fortschr. Phys.}, vol. 48, no. 9--11, pp. 771--783, 2000,
doi: 10.1002/1521-3978(200009)48:9/11<771::AID-PROP771>3.0.CO;2-E.

\bibitem{preskill}
J. Preskill, "Quantum computing in the NISQ era and beyond,"
\emph{Quantum}, vol. 2, Art. no. 79, 2018,
doi: 10.22331/q-2018-08-06-79.

\bibitem{hangleiter}
D. Hangleiter and J. Eisert, "Computational advantage of quantum random
sampling," \emph{Rev. Mod. Phys.}, vol. 95, no. 3, Art. no. 035001, 2023,
doi: 10.1103/RevModPhys.95.035001.

\bibitem{zhong}
H.-S. Zhong \emph{et al.}, "Quantum computational advantage using photons,"
\emph{Science}, vol. 370, no. 6523, pp. 1460--1463, 2020,
doi: 10.1126/science.abe8770.

\bibitem{fowler}
A. G. Fowler, M. Mariantoni, J. M. Martinis, and A. N. Cleland, "Surface
codes: Towards practical large-scale quantum computation," \emph{Phys.
Rev. A}, vol. 86, Art. no. 032324, 2012,
doi: 10.1103/PhysRevA.86.032324.

\bibitem{bravyi}
S. Bravyi \emph{et al.}, "High-threshold and low-overhead fault-tolerant
quantum memory," \emph{Nature}, vol. 627, no. 8005, pp. 778--782, 2024,
doi: 10.1038/s41586-024-07107-7.

\bibitem{cohen}
L. Z. Cohen, I. H. Kim, S. D. Bartlett, and B. J. Brown, "Low-overhead
fault-tolerant quantum computing using long-range connectivity,"
\emph{Sci. Adv.}, vol. 8, no. 20, Art. no. eabn1717, 2022,
doi: 10.1126/sciadv.abn1717.
\end{thebibliography}

\end{document}
